\documentclass[10pt,a4paper]{amsart}
\usepackage[margin=19mm]{geometry}
\usepackage{amsmath,amssymb,mathtools}
\usepackage[scr=rsfs,cal=euler]{mathalfa}
\usepackage{xfrac}
\usepackage[unicode,hidelinks,bookmarks=false]{hyperref}
\newcommand{\ie}{i.e.\ }
\newcommand{\cf}{cf.\ }
\newcommand{\resp}{respectively\ }
\newcommand{\Subsection}{\subsection}
\providecommand{\nobreakdash}{\nobreak\hbox{-}}
\setlength{\emergencystretch}{2em}
\multlinegap0pt
\usepackage{braket}
\usepackage{amssymb}
\usepackage{algorithm}\usepackage{algpseudocode}
\newtheorem{theorem}{Theorem}
\newtheorem{remark}{Remark}
\DeclareMathOperator{\Tr}{Tr}
\title{Bridging Classical and Quantum:\ Group-Theoretic Approach to Quantum Circuit Simulation}
\author{Daksh Shami}
\date{Source: arXiv:2407.19575v3 (CC BY 4.0)}
\begin{document}
\maketitle

\noindent\emph{Source and licence.} Bridging Classical and Quantum:\ Group-Theoretic Approach to Quantum Circuit Simulation. Version arXiv:2407.19575v3 (CC BY 4.0), distributed by arXiv under the Creative Commons Attribution 4.0 International licence, http://creativecommons.org/licenses/by/4.0/, as recorded in the arXiv licence metadata for this exact version and shown on the abstract page https://arxiv.org/abs/2407.19575v3. Full licence notice: \texttt{licenses/arxiv-2407.19575-CC-BY-4.0.txt}.\par\smallskip

\noindent\textit{Editorial scope.} Counted unit: the article's theorem theorem (source0.tex lines 240--248). The article's complete original proof of it is delivered with it (source0.tex lines 249--313, one \texttt{proof} environment of 65 lines). No mathematics is written, repaired or re-derived: the statement and the proof are the article's own lines, in the article's own notation, and no retained line is substituted or re-flowed.  Retained uncounted as the source-local context the proof needs: lines 48--60.  Nothing was omitted from any retained span, so the package carries no omission record.  No retained line contains a citation, so the wrapper carries no bibliography and no bibliography entry.  No retained line contains a figure environment, an image-inclusion command or drawing code, so no image is delivered and the package declares no asset dependency.  Editorial formatting only, in addition: an AMS \texttt{amsart} preamble that restates the article's own declarations and macros used by the retained lines, namely the environments \texttt{theorem}, \texttt{remark} on separate counters, together with the standard packages those lines need.  The retained environments print in this document as Theorem 1, Remark 1 in order of appearance, whereas the article prints the same items with the numbering of its own full document.  The complete native source of the article as distributed in the arXiv e-print for this exact version is bundled unchanged as \texttt{sources/163008/source0.tex}.
\begin{theorem}[Character Function Decomposition]
\label{thm:char_decomp}
Let $G$ be a finite group and $\mathbb{C}[G]$ its complex group algebra. For any element $u \in \mathbb{C}[G]$, it can be decomposed into a sum of character-weighted representations as follows:
\begin{equation}
u = \sum_{i=1}^k \frac{\chi_i(u)}{d_i} \sum_{g \in G} \chi_i(g^{-1}) \rho_i(g),
\end{equation}
where:
\begin{itemize}
    \item $k$ is the number of irreducible representations of $G$,
    \item $\chi_i$ is the character of the $i$-th irreducible representation $\rho_i$,
    \item $d_i$ is the dimension of $\rho_i$.
\end{itemize}
\end{theorem}

\begin{theorem}[Generalized Gottesman-Knill]
Let $G$ be a finite group with a set of generators ${g_1, \ldots, g_k}$. If:
\begin{enumerate}
\item The irreducible representations of $G$ can be efficiently computed,
\item The character values $\chi_i(g)$ can be efficiently computed for all $g \in G$ and all irreducible representations $i$,
\item The number of irreducible representations grows polynomially with the number of qubits,
\end{enumerate}
Then quantum circuits composed of gates from $G$ can be efficiently simulated classically using the character function decomposition method.
\end{theorem}

\begin{proof}
Let $U$ be a quantum circuit composed of $m$ gates from $G$, where $m$ is polynomial in the number of qubits $n$. We can express $U$ as:
\begin{equation}
U = g_{i_1} g_{i_2} \cdots g_{i_m}
\end{equation}
where each $g_{i_j}$ is a generator of $G$. By Theorem \ref{thm:char_decomp}, we can decompose $U$ as:
\begin{equation}
U = \sum_{i=1}^k \frac{\chi_i(U)}{d_i} \sum_{g \in G} \chi_i(g^{-1}) \rho_i(g)
\end{equation}
where $k$ is the number of irreducible representations of $G$, $\chi_i$ is the character of the $i$-th irreducible representation $\rho_i$, and $d_i$ is the dimension of $\rho_i$.
To simulate this circuit efficiently, we need to show that:
\begin{enumerate}
\item We can compute $\chi_i(U)$ efficiently for all $i$.
\item We can evaluate the sum $\sum_{g \in G} \chi_i(g^{-1}) \rho_i(g)$ efficiently for all $i$.
\item We can efficiently compute the expectation value of an observable $O$ for any input state $\ket{\psi}$.
\end{enumerate}

\noindent

Computing $\chi_i(U)$:

By the properties of characters, we have:
\begin{equation}
\chi_i(U) = \chi_i(g_{i_1} g_{i_2} \cdots g_{i_m}) = \prod_{j=1}^m \chi_i(g_{i_j})
\end{equation}
By assumption 2, each $\chi_i(g_{i_j})$ can be computed efficiently. Since $m$ is polynomial in $n$, the product can be computed efficiently.

\begin{remark}
\label{rmk:char_mult}
Equation~(12) relies on the multiplicativity of characters, 
$\chi_i(gh) = \chi_i(g)\chi_i(h)$, which holds if and only if 
$\rho_i$ is one-dimensional. For irreducible representations of 
dimension $d_i > 1$, we have \[\chi_i(gh) = \Tr(\rho_i(g)\rho_i(h)) 
\neq \Tr(\rho_i(g))\Tr(\rho_i(h)) = \chi_i(g)\chi_i(h)\] in general. 
Consequently, the efficient simulation argument in Theorems~4 and~5 
applies as stated to groups whose irreducible representations are all 
one-dimensional i.e. abelian groups. The structure of the 
deviation $\chi(gh) - \chi(g)\chi(h)$ for the non-abelian case, and 
whether it admits an efficient approximation scheme, is an interesting 
open direction.
\end{remark}

\noindent
2. Evaluating $\sum_{g \in G} \chi_i(g^{-1}) \rho_i(g)$:
Let $S_i = \sum_{g \in G} \chi_i(g^{-1}) \rho_i(g)$. We don't need to compute this sum explicitly. Instead, we can use the fact that $S_i$ is a scalar multiple of the identity matrix in the $i$-th irreducible representation:
\begin{equation}
S_i = \frac{|G|}{d_i} I_{d_i}
\end{equation}
This follows from Schur's lemma and the orthogonality relations of characters. We can compute $|G|/d_i$ efficiently since $d_i \leq \sqrt{|G|}$ for all $i$.
\noindent
3. Computing expectation values:
For an observable $O$ and input state $\ket{\psi}$, we need to compute $\bra{\psi}U^\dagger O U\ket{\psi}$ efficiently. Using the decomposition of $U$:
\[
\bra{\psi}U^\dagger O U\ket{\psi} = \sum_{i,j=1}^k \frac{\overline{\chi_i(U)}\chi_j(U)}{d_i d_j} \bra{\psi}\left(\sum_{g \in G} \chi_i(g) \rho_i(g)^\dagger\right) O \left(\sum_{h \in G} \chi_j(h^{-1}) \rho_j(h)\right)\ket{\psi} \]\[
= \sum_{i,j=1}^k \frac{\overline{\chi_i(U)}\chi_j(U)}{d_i d_j} \frac{|G|^2}{d_i d_j} \bra{\psi}O\ket{\psi} 
= |G|^2 \bra{\psi}O\ket{\psi} \sum_{i=1}^k \frac{|\chi_i(U)|^2}{d_i^2}
\]
The last step uses the fact that $\sum_i (d_i^2 / |G|) = 1$. We can compute this efficiently because:
\begin{enumerate}
\item $\chi_i(U)$ can be computed efficiently (from step 1).
\item The number of irreducible representations $k$ is polynomial in $n$ (by assumption 3).
\item $\bra{\psi}O\ket{\psi}$ can be computed efficiently for typical observables and input states.
\end{enumerate}
Therefore, we can efficiently compute expectation values for any observable $O$ and input state $\ket{\psi}$, which allows for efficient classical simulation of the quantum circuit.
\end{proof}

\end{document}
